%============================================================
% CH2 · P3 · PROBLEM 2 — COUNT DESCENDANTS (AtCoder ABC202 E)
%============================================================
\subsection{Problem 2.3.2: Count Descendants}
\problemlabel{2.3.2}{ssec:ch2-count-descendants}

\problemcard{2.3.2}{Count Descendants}{AtCoder ABC202 E}{https://atcoder.jp/contests/abc202/tasks/abc202_e}{\exmedium\ (AtCoder ABC 500 points)}{Entry times per depth, sorted for free; count those inside $[\tin{U}, \tout{U}]$ with two binary searches --- or offline, a counter difference between entering and leaving $U$}

\tikzset{
    g2p3c-v1/.style=gv, g2p3c-v2/.style=gv, g2p3c-v3/.style=gv, g2p3c-v4/.style=gv,
    g2p3c-v5/.style=gv, g2p3c-v6/.style=gv, g2p3c-v7/.style=gv,
    g2p3c-T/.pic={
        \node[g2p3c-v1] (c1) at (2.4,0) {1};
        \node[g2p3c-v2] (c2) at (1.2,-1.2) {2};
        \node[g2p3c-v3] (c3) at (3.6,-1.2) {3};
        \node[g2p3c-v4] (c4) at (0.2,-2.4) {4};
        \node[g2p3c-v5] (c5) at (1.2,-2.4) {5};
        \node[g2p3c-v7] (c7) at (2.2,-2.4) {7};
        \node[g2p3c-v6] (c6) at (0.2,-3.6) {6};
        \foreach \a/\b in {1/2, 1/3, 2/4, 2/5, 2/7, 4/6} \draw[gtree, -{Stealth[length=2.2mm]}] (c\a) -- (c\b);
    }
}

\subsubsection*{1. Problem statement}

A tree with $N$ vertices is rooted at vertex 1; vertex $i \ge 2$ has parent $P_i < i$. Answer $Q$ queries $(U_i, D_i)$: how many vertices $u$ have $U_i$ on their path to the root (endpoints included) \emph{and} exactly $D_i$ edges between $u$ and the root?

\textbf{Constraints:} $2 \le N \le 2\cdot10^5$, $1 \le P_i < i$, $1 \le Q \le 2\cdot10^5$, $1 \le U_i \le N$, $0 \le D_i \le N-1$.

\begin{table}[H]
\centering\small
\begin{tabular}{@{}lll@{}}
\toprule
\textbf{Input} & \textbf{Official output} & \textbf{Our output} \\
\midrule
\texttt{7 / 1 1 2 2 4 2 / 4 / 1 2 / 7 2 / 4 1 / 5 5} & \texttt{3 1 0 0} & \texttt{3 1 0 0} \\
\bottomrule
\end{tabular}
\end{table}

\begin{diagrambox}
\begin{tikzpicture}[x=1cm, y=0.9cm]
    \begin{scope}
        \pic {g2p3c-T};
        \foreach \d in {0,1,2,3} \node[font=\scriptsize, text=darkgray, anchor=east] at (-0.4,-1.2*\d) {depth \d};
        \node[font=\scriptsize\bfseries] at (2.0,0.8) {the sample tree};
    \end{scope}
    \begin{scope}[xshift=6.2cm, g2p3c-v1/.style=gvsrc, g2p3c-v4/.style=gvdst, g2p3c-v5/.style=gvdst, g2p3c-v7/.style=gvdst]
        \pic {g2p3c-T};
        \draw[accentorange, dashed, rounded corners=3pt] (-0.4,-2.0) rectangle (2.7,-2.8);
        \node[font=\scriptsize\bfseries] at (2.0,0.8) {query $(1, 2)$: answer 3};
        \node[note, anchor=west, text width=3.6cm, align=left] at (3.2,-2.6) {every vertex of depth 2 lies below the root: $4, 5, 7$};
    \end{scope}
\end{tikzpicture}
\figcap{``$U$ is on the path from $u$ to the root'' means ``$u$ is a descendant of $U$ (or $U$ itself)''. Query $(7, 2)$ counts only 7; query $(4, 1)$ asks for depth 1 below a vertex of depth 2: 0; depth 5 does not exist: 0.}
\end{diagrambox}

\subsubsection*{2. The graph model}

\begin{modelbox}[Graph Model --- Count Descendants]
\renewcommand{\arraystretch}{1.35}
\begin{tabularx}{\linewidth}{@{}>{\bfseries}l X@{}}
Vertices & the $N$ tree vertices, each with a depth $\mathrm{dep}[v]$ \\
Edges & $P_i \to i$, rooted at 1 \\
Weights & none \\
Question & $\#\{v : v \text{ in the subtree of } U,\ \mathrm{dep}[v] = D\}$ --- a subtree condition and a group condition \\
Size & $N, Q \le 2\cdot10^5$: \bigO{N} or \bigO{N \log N} preprocessing, \bigO{\log N} per query; depth up to $N-1$ \\
\end{tabularx}
\end{modelbox}

\begin{diagrambox}
\begin{tikzpicture}[x=0.75cm, y=0.5cm]
    \draw[->, darkgray] (-0.3,0) -- (14.4,0) node[right, font=\scriptsize] {time};
    \foreach \t in {0,...,13} {\draw[darkgray] (\t,0.12) -- (\t,-0.12); \node[font=\tiny, below] at (\t,-0.12) {\t};}
    \foreach \v/\a/\b/\l/\c in {1/0/13/4/primaryblue, 2/1/10/3/successgreen, 3/11/12/3/successgreen,
                                4/2/5/2/accentorange, 5/6/7/2/accentorange, 7/8/9/2/accentorange, 6/3/4/1/softpurple}{
        \draw[\c, line width=2.2pt] (\a,\l*1.0) -- (\b,\l*1.0);
        \draw[\c, line width=1pt] (\a,\l*1.0-0.25) -- (\a,\l*1.0+0.25);
        \draw[\c, line width=1pt] (\b,\l*1.0-0.25) -- (\b,\l*1.0+0.25);
        \node[font=\scriptsize\bfseries, fill=white, inner sep=1pt, text=\c] at ({(\a+\b)/2},\l*1.0+0.45) {\v};
    }
    \foreach \d/\l in {0/4, 1/3, 2/2, 3/1} \node[font=\scriptsize, text=darkgray, anchor=east] at (-0.4,\l*1.0) {dep \d};
\end{tikzpicture}
\figcap{The model in time: intervals $[\tin{v}, \tout{v}]$ of the sample, one row per depth. A query is ``how many intervals of row $D$ start inside the interval of $U$''.}
\end{diagrambox}

\subsubsection*{3. Key observations}

\begin{itemize}[leftmargin=*]
    \item \textbf{Path to the root through $U$ $=$ descendant of $U$.} Then the parenthesis theorem turns the subtree condition into $\tin{U} \le \tin{v} \le \tout{U}$.
    \item \textbf{Group by depth.} For each depth $d$ keep the list $L_d$ of entry times of the depth-$d$ vertices. The DFS appends to $L_d$ in time order, so every list is sorted without sorting.
    \item \textbf{Count $=$ two binary searches.} $\texttt{upper\_bound}(L_D, \tout{U}) - \texttt{lower\_bound}(L_D, \tin{U})$. A depth below $U$'s own depth gives 0 automatically (no shallower vertex is inside the window), and so does a depth that does not occur (empty list).
    \item \textbf{Why not per-vertex tables?} Storing ``count per depth'' for every vertex is \bigO{N^2} memory on a path.
    \item \textbf{Depth up to $2\cdot10^5$}: iterative DFS.
\end{itemize}

\subsubsection*{4. Approaches}

\begin{bruteforcebox}
\textbf{Idea: test every vertex.} For a query, take each vertex of depth $D$ and climb to the root; count it if $U$ is met.

\begin{lstlisting}[style=cpp]
#include <bits/stdc++.h>
using namespace std;

// Brute force: for every query, look at every vertex v of depth D and walk from v up to the root;
// count v if U is met on the way.
int main() {
    int n;
    cin >> n;
    vector<int> p(n + 1, 0), dep(n + 1, 0);
    for (int i = 2; i <= n; i++) {
        cin >> p[i];
        dep[i] = dep[p[i]] + 1;                  // p_i < i: the parent's depth is known
    }
    int q;
    cin >> q;
    while (q--) {
        int u, d, ans = 0;
        cin >> u >> d;
        for (int v = 1; v <= n; v++) {
            if (dep[v] != d) continue;
            for (int x = v; x != 0; x = p[x])
                if (x == u) { ans++; break; }
        }
        cout << ans << "\n";
    }
}
\end{lstlisting}

\begin{complexitybox}
\textbf{Time:} \bigO{N \cdot \text{depth}} per query --- on a path with $D = N-1$ that is $N$ steps per query, $4\cdot10^{10}$ in total. \textbf{Space:} \bigO{N}. \textbf{Verdict:} \tle{}; reference only.
\end{complexitybox}
\end{bruteforcebox}

\begin{algorithmbox}[Approach 2: sorted entry times per depth + binary search]
\textbf{Idea.} One iterative DFS records $\tin{}$, $\tout{}$, $\mathrm{dep}$ and appends $\tin{v}$ to \texttt{at[dep[v]]} at the moment $v$ is entered. Each query is two binary searches in one list.

\begin{lstlisting}[style=cpp]
#include <bits/stdc++.h>
using namespace std;

// Count Descendants (ABC202 E): count v at depth D with U on the path root..v,
// i.e. descendants of U (U included) at depth D. Iterative DFS with one clock;
// at[d] = entry times of the depth-d vertices, sorted because they are appended in time order.
int main() {
    ios::sync_with_stdio(false);
    cin.tie(nullptr);
    int n;
    cin >> n;
    vector<vector<int>> ch(n + 1);
    for (int i = 2; i <= n; i++) {
        int p;
        cin >> p;
        ch[p].push_back(i);
    }
    vector<int> tin(n + 1), tout(n + 1), dep(n + 1, 0);
    vector<size_t> it(n + 1, 0);
    vector<vector<int>> at(n);
    int timer = 0;
    vector<int> st(1, 1);
    tin[1] = timer++; at[0].push_back(tin[1]);
    while (!st.empty()) {
        int u = st.back();
        if (it[u] < ch[u].size()) {
            int v = ch[u][it[u]++];
            dep[v] = dep[u] + 1;
            tin[v] = timer++;
            at[dep[v]].push_back(tin[v]);
            st.push_back(v);
        } else {
            tout[u] = timer++;
            st.pop_back();
        }
    }
    int q;
    cin >> q;
    while (q--) {
        int u, d;
        cin >> u >> d;
        const vector<int>& L = at[d];             // d <= n - 1, so the index is valid
        int ans = upper_bound(L.begin(), L.end(), tout[u]) - lower_bound(L.begin(), L.end(), tin[u]);
        cout << ans << "\n";
    }
}
\end{lstlisting}

\begin{complexitybox}
\textbf{Time:} \bigO{N} for the DFS (the lists hold $N$ numbers in total), \bigO{\log N} per query: \bigO{N + Q \log N}. \textbf{Space:} \bigO{N}. \textbf{Verdict:} \ac{}-ready; on the path $P_i = i-1$ with $2\cdot10^5$ random queries the run takes under a second here.
\end{complexitybox}
\end{algorithmbox}

\begin{invariantbox}[Why the count is right]
\textbf{Lists are sorted.} \texttt{timer} only grows, and $\tin{v}$ is appended to \texttt{at[dep[v]]} at the moment it is assigned; so each list is strictly increasing.

\textbf{The window selects the subtree.} By the parenthesis theorem, $v$ is a descendant of $U$ (or $U$ itself) iff $\tin{U} \le \tin{v} \le \tout{U}$. Restricted to the list of depth $D$, the vertices that qualify are exactly the entries in $[\tin{U}, \tout{U}]$; in a sorted list they form the index range $[\texttt{lower\_bound}(\tin{U}), \texttt{upper\_bound}(\tout{U}))$, whose length is the answer. Depths smaller than $\mathrm{dep}[U]$ have no entry in the window (no ancestor of a vertex can be its descendant), and missing depths are empty lists. \hfill$\blacksquare$
\end{invariantbox}

\begin{alternativebox}[Approach 3: offline --- counter difference between entering and leaving $U$]
Collect the queries at their vertex $U$. During the DFS keep \texttt{cnt[d]} $=$ number of depth-$d$ vertices entered so far. Everything entered after ``enter $U$'' and before ``leave $U$'' is in $U$'s subtree, so
\[
    \text{answer}(U, D) = \texttt{cnt}[D]\big|_{\text{leaving } U} - \texttt{cnt}[D]\big|_{\text{entering } U,\ \text{before counting } U}.
\]
No sorting, no binary search: \bigO{N + Q}.

\begin{lstlisting}[style=cpp]
#include <bits/stdc++.h>
using namespace std;

// Offline alternative: cnt[d] = number of depth-d vertices entered so far.
// Everything entered between "enter U" and "leave U" is in U's subtree, so
// answer(U, D) = cnt[D] when leaving U  -  cnt[D] when entering U.
int main() {
    ios::sync_with_stdio(false);
    cin.tie(nullptr);
    int n;
    cin >> n;
    vector<vector<int>> ch(n + 1);
    for (int i = 2; i <= n; i++) {
        int p;
        cin >> p;
        ch[p].push_back(i);
    }
    int q;
    cin >> q;
    vector<vector<pair<int, int>>> qs(n + 1);     // qs[u] = (D, query index)
    vector<long long> ans(q, 0);
    for (int i = 0; i < q; i++) {
        int u, d;
        cin >> u >> d;
        qs[u].push_back({d, i});
    }
    vector<int> cnt(n, 0), dep(n + 1, 0);
    vector<size_t> it(n + 1, 0);
    vector<int> st;
    auto enter = [&](int v) {
        for (auto& pr : qs[v]) ans[pr.second] -= cnt[pr.first];   // before v itself is counted
        cnt[dep[v]]++;
        st.push_back(v);
    };
    enter(1);
    while (!st.empty()) {
        int u = st.back();
        if (it[u] < ch[u].size()) {
            int v = ch[u][it[u]++];
            dep[v] = dep[u] + 1;
            enter(v);
        } else {
            for (auto& pr : qs[u]) ans[pr.second] += cnt[pr.first];
            st.pop_back();
        }
    }
    for (int i = 0; i < q; i++) cout << ans[i] << "\n";
}
\end{lstlisting}

\begin{complexitybox}
\textbf{Time:} \bigO{N + Q}. \textbf{Space:} \bigO{N + Q}. \textbf{Verdict:} \ac{}-ready, identical output to Approach~2 on the $2\cdot10^5$ path test. It needs all queries in advance; Approach~2 is \emph{online} (it could answer a query before reading the next). The difference trick generalises to any group key (colour, letter) and is the seed of the ``DSU on tree'' / small-to-large techniques of Chapter~15.
\end{complexitybox}
\end{alternativebox}

\begin{diagrambox}
\begin{tikzpicture}[x=0.78cm, y=0.8cm]
    \foreach \t/\ev in {0/{in 1}, 1/{in 2}, 2/{in 4}, 3/{in 6}, 4/{out 6}, 5/{out 4}, 6/{in 5}, 7/{out 5}, 8/{in 7}, 9/{out 7}, 10/{out 2}, 11/{in 3}, 12/{out 3}, 13/{out 1}}{
        \node[font=\tiny, rotate=60, anchor=west] at (\t,0.25) {\ev};
        \draw[darkgray] (\t,0.12) -- (\t,-0.12);
    }
    \draw[darkgray] (-0.3,0) -- (13.3,0);
    \node[font=\scriptsize, anchor=east] at (-0.4,-0.7) {cnt[2] after};
    \foreach \t/\c in {0/0, 1/0, 2/1, 3/1, 4/1, 5/1, 6/2, 7/2, 8/3, 9/3, 10/3, 11/3, 12/3, 13/3}
        \node[font=\scriptsize\ttfamily] at (\t,-0.7) {\c};
    \draw[accentorange, line width=1.2pt, -{Stealth[length=2mm]}] (7,-1.2) -- (7,-0.95);
    \draw[accentorange, line width=1.2pt, -{Stealth[length=2mm]}] (9,-1.2) -- (9,-0.95);
    \node[font=\scriptsize, text=accentorange, anchor=north, text width=8cm, align=center] at (8,-1.25) {query $(7, 2)$: on entering 7, before 7 itself is counted, \texttt{cnt[2]} is still 2 (the value after event 7); on leaving 7 it is 3: $3 - 2 = 1$};
\end{tikzpicture}
\figcap{Approach 3 on the sample: the depth-2 counter along the 14 events of the DFS. Query $(1,2)$ reads $0$ and $3$; query $(4,1)$ reads \texttt{cnt[1]} $= 1$ on both sides: $0$.}
\end{diagrambox}

\begin{tricksbox}[Approach 4: Implementation Refinements]
\begin{itemize}[leftmargin=*]
    \item \textbf{Size the depth table by $N$.} $D \le N - 1$, so \texttt{at} with $N$ lists can be indexed directly; with \texttt{at} sized by the maximum depth, guard the index.
    \item \textbf{Pre-order positions work too.} With $\mathrm{pos}$ and $\mathrm{sz}$ the window is $[\mathrm{pos}[U], \mathrm{pos}[U]+\mathrm{sz}[U]-1]$ --- then use \texttt{upper\_bound} on the right end (Mistake 3 of the pattern).
    \item \textbf{No DFS needed for depths} ($P_i < i$), but the entry times need the DFS; the recursion depth can be $N$: iterative.
    \item \textbf{Memory.} $N$ entries in all lists together plus the children lists: a few MB.
\end{itemize}
\end{tricksbox}

\subsubsection*{5. Trace}

\begin{tracebox}[{Sample: children 1: 2 3,\ 2: 4 5 7,\ 4: 6}]
\begin{center}
\begin{tikzpicture}[x=0.72cm, y=0.72cm]
    \tikzset{gv/.append style={minimum size=6mm, font=\scriptsize\bfseries}}
    \begin{scope}[g2p3c-v1/.style=gvseen, g2p3c-v2/.style=gvseen, g2p3c-v4/.style=gvseen, g2p3c-v6/.style=gvcur]
        \pic {g2p3c-T};
        \node[font=\scriptsize\bfseries] at (2.0,0.9) {t = 3: enter 6};
        \node[font=\scriptsize\ttfamily] at (2.0,-4.5) {stack 1 2 4 6};
    \end{scope}
    \begin{scope}[xshift=4.4cm, g2p3c-v1/.style=gvseen, g2p3c-v2/.style=gvseen, g2p3c-v4/.style=gvdone, g2p3c-v6/.style=gvdone, g2p3c-v5/.style=gvdone, g2p3c-v7/.style=gvcur]
        \pic {g2p3c-T};
        \node[font=\scriptsize\bfseries] at (2.0,0.9) {t = 8: enter 7};
        \node[font=\scriptsize\ttfamily] at (2.0,-4.5) {stack 1 2 7};
    \end{scope}
    \begin{scope}[xshift=8.8cm, g2p3c-v1/.style=gvcur, g2p3c-v2/.style=gvdone, g2p3c-v3/.style=gvdone, g2p3c-v4/.style=gvdone, g2p3c-v6/.style=gvdone, g2p3c-v5/.style=gvdone, g2p3c-v7/.style=gvdone]
        \pic {g2p3c-T};
        \node[font=\scriptsize\bfseries] at (2.0,0.9) {t = 13: leave 1};
        \node[font=\scriptsize\ttfamily] at (2.0,-4.5) {stack (empty)};
    \end{scope}
\end{tikzpicture}
\end{center}
\begin{center}\small
\begin{tabular}{@{}lccccccc@{}}
\toprule
vertex & 1 & 2 & 3 & 4 & 5 & 6 & 7 \\
\midrule
$\tin{}$ & 0 & 1 & 11 & 2 & 6 & 3 & 8 \\
$\tout{}$ & 13 & 10 & 12 & 5 & 7 & 4 & 9 \\
depth & 0 & 1 & 1 & 2 & 2 & 3 & 2 \\
\bottomrule
\end{tabular}
\qquad
\begin{tabular}{@{}ll@{}}
\toprule
\textbf{depth} & \textbf{\texttt{at[d]}} \\
\midrule
0 & 0 \\
1 & 1, 11 \\
2 & 2, 6, 8 \\
3 & 3 \\
\bottomrule
\end{tabular}
\end{center}
\begin{center}\small
\begin{tabular}{@{}ccll@{}}
\toprule
\textbf{query} & \textbf{window} & \textbf{list} & \textbf{answer} \\
\midrule
$(1,2)$ & $[0, 13]$ & $2, 6, 8$: all inside & 3 \\
$(7,2)$ & $[8, 9]$ & only 8 inside & 1 \\
$(4,1)$ & $[2, 5]$ & $1, 11$: none inside & 0 \\
$(5,5)$ & $[6, 7]$ & \texttt{at[5]} empty & 0 \\
\bottomrule
\end{tabular}
\end{center}
\end{tracebox}

\subsubsection*{6. Test suite}

\begin{testbox}
\centering\small
\begin{tabular}{@{}p{6.4cm}lp{5cm}@{}}
\toprule
\textbf{Input} & \textbf{Our output} & \textbf{Why} \\
\midrule
\texttt{2 / 1 / 3 / 1 0 / 1 1 / 2 0} & \texttt{1 1 0} & $D = 0$ counts the root itself; depth above $U$ \\
\texttt{5 / 1 2 3 4 / 3 / 1 4 / 3 2 / 5 4} & \texttt{1 1 1} & a path: $U$ counts itself when $D = \mathrm{dep}[U]$ \\
\texttt{5 / 1 1 1 1 / 3 / 1 1 / 2 1 / 1 4} & \texttt{4 1 0} & a star; a missing depth \\
\texttt{7 / 1 1 2 2 4 2 / 3 / 2 2 / 3 2 / 6 3} & \texttt{3 0 1} & the sample tree, other queries \\
path $P_i = i-1$, $N = Q = 2\cdot10^5$ & (same for 2 and 3) & depth $2\cdot10^5$; iterative only \\
\bottomrule
\end{tabular}
\end{testbox}
All outputs come from running Approach~2. Approaches 1--3 reproduce the official sample and agree on $1000$ random trees with $N \le 10$ and $Q \le 10$ random queries (paths and stars included).

\subsubsection*{7. Common mistakes}

\begin{warningbox}[Mistake 1 --- forgetting $U$ itself]
``Endpoints included'': with $D = \mathrm{dep}[U]$ the answer is 1 ($U$ alone). The window $[\tin{U}, \tout{U}]$ includes $\tin{U}$, so the formula handles it; a ``strict descendants'' implementation answers \texttt{1 0 0} instead of \texttt{1 1 1} on the second test.
\end{warningbox}

\begin{warningbox}[Mistake 2 --- sorting what is already sorted, or sorting by vertex number]
The lists must be sorted by \emph{time}. Building them in a loop over $v = 1..N$ (``\texttt{at[dep[v]].push\_back(tin[v])}'' after the DFS) gives lists ordered by vertex number, and binary search on them is wrong unless you sort. Appending during the DFS needs no sort.
\end{warningbox}

\begin{warningbox}[Mistake 3 --- the offline version counts $U$ before subtracting]
In Approach 3, \texttt{ans -= cnt[D]} must happen \emph{before} \texttt{cnt[dep[U]]++}; otherwise $U$ itself is never counted, and $(7, 2)$ of the sample becomes 0.
\end{warningbox}

\begin{warningbox}[Mistake 4 --- recursion]
$P_i = i - 1$ is allowed: a recursive DFS goes $2\cdot10^5$ deep and crashes on a 2\,MB stack.
\end{warningbox}

\subsubsection*{8. Takeaways}

\begin{itemize}[leftmargin=*]
    \item ``$U$ on the path from $u$ to the root'' is the subtree of $U$: an interval of times.
    \item Subtree $\times$ group $=$ one sorted list of times per group, count inside the window by two binary searches.
    \item Offline: a counter read on entering and leaving is the same count without searching.
    \item Lists filled during the DFS are sorted by time for free.
\end{itemize}

\begin{variantsbox}
\variant{Depth relative to $U$}{\varnone}
  {count the descendants exactly $K$ levels below $U$.}
  {the same query with $D = \mathrm{dep}[U] + K$.}
  {one addition before the lookup; $D > N-1$ gives 0.}
\variant{Count by colour instead of depth}{\varnone}
  {every vertex has a colour; count descendants of $U$ with colour $c$.}
  {group the vertices by colour instead of depth.}
  {lists of entry times per colour, the same two binary searches; or the offline counter per colour (checked against the brute force on random coloured trees).}
\variant{The $k$-th descendant at depth $D$}{\varnone}
  {in DFS order, the $k$-th vertex of depth $D$ below $U$.}
  {the qualifying entries are a contiguous run of $L_D$.}
  {index \texttt{lower\_bound}$(\tin{U}) + k - 1$, valid if it is still $\le \tout{U}$; store the vertex next to its time.}
\variant{Ancestor test only}{\varnone}
  {just ``is $U$ on the root path of $v$''.}
  {no group at all.}
  {two comparisons of times (Exercise E3.1).}
\end{variantsbox}
